\documentclass[letterpaper]{article}

\topmargin = 0.8cm
\oddsidemargin=0.60in % leftmargin is 1 inch
\textwidth=5.5in
\linespread{1.1} % Line spacing
\usepackage{booktabs}
\usepackage{multirow}
\usepackage{amsfonts,amssymb,latexsym,amscd}
\usepackage{bm}
\usepackage{amsmath,mathtools}
\usepackage{soul} 
\usepackage{subfigure}
\usepackage[textsize=tiny,textwidth=1.4cm]{todonotes}
\usepackage{microtype}
\usepackage{graphicx}
\usepackage{epstopdf}
\usepackage{pdfpages}
\usepackage{longtable}
\usepackage{fancyhdr}
\usepackage{color}
\usepackage{enumerate}
\usepackage{lipsum}
\usepackage{alltt}
\usepackage{tikz}
\usepackage{url}
\usetikzlibrary{calc,patterns,decorations.pathmorphing,decorations.markings,decorations.pathreplacing}
\usepackage{cancel}
\usepackage{listings}
\usepackage{xcolor}
\usepackage[linkcolor=black,colorlinks=true,urlcolor=blue]{hyperref}
\usepackage{wrapfig}
\usepackage{lastpage}
\usepackage[yyyymmdd,hhmmss]{datetime}
\usepackage[algoruled,linesnumbered,vlined,titlenotnumbered]{algorithm2e}


\cfoot{}
\cfoot{\thepage\ of \pageref{LastPage}}
%\rfoot{[\today]}


\oddsidemargin=-0.25in % leftmargin is 1 inch
\voffset=-0.9in % leftmargin is 1 inch
\textwidth=7in
\textheight=9.25in
\headsep = 36pt

\newcommand\independent{\protect\mathpalette{\protect\independenT}{\perp}}
\def\independenT#1#2{\mathrel{\rlap{$#1#2$}\mkern2mu{#1#2}}}

\newcommand{\courseS}{Hierarchical Softmax Calibration}
\newcommand{\deptS}{Polytechnique Montréal}
\newcommand{\profS}{J.A. Goulet, L.H. Nguyen \& M. Florensa-Motilla}
\newcommand{\dateS}{\today}
\newcommand{\titleS}{~}


\author{\courseS \\ \deptS\\ Enseignant: \profS\\[6pt] \dateS}
\date{}

\pagestyle{fancy}
\fancyhead{}
\fancyhead[c]{ \sc \courseS \hfill \profS\\ \raggedright \deptS \hfill \dateS  }
\fancyfoot{}
\fancyfoot[l]{\tiny (Last update : \today, \currenttime)\\[6pt]}
\fancyfoot[c]{\thepage/\pageref{LastPage}}

\newcommand{\titleshort}{\titleS}

\begin{document}
\title{\titleS }

\section{Context}\label{S:Context}
In the hierarchical binary decomposition employed by TAGI for classification (Goulet, Nguyen \& Amiri, JMLR 2021, Appendix~C), a class $\mathcal C=(\mathcal C_1,\dots,\mathcal C_{\mathtt H})\in\{0,1\}^{\mathtt H}$, $\mathtt H=\lceil\log_2\mathtt K\rceil$, is identified with a path in a binary tree. Each node $h$ on the path is associated with an output-layer hidden state $Z_h\sim\mathcal N(\mu_h,\sigma_h^2)$ and a branch sign $s_h=(-1)^{\mathcal C_h}$, and the class probability is
\begin{equation}
p(\mathcal C\mid\mathcal D)=\prod_{h=1}^{\mathtt H}\Phi\!\left(\frac{s_h\,\mu_h}{\sqrt{\alpha^2+\sigma_h^2}}\right),
\label{E:tagi}
\end{equation}
where $\Phi$ is the standard normal CDF and $\alpha$ a fixed scaling hyperparameter. Each factor of \eqref{E:tagi} is the expected value $\mathbb E[\Phi(s_hZ_h/\alpha)]$: the uncertainty on $Z_h$ is marginalized, but the class probability itself carries no variance, and the scale $\alpha$ has to be tuned by hand.

This note does two things.
\begin{enumerate}[(i)]
\item Each factor $\Phi(\cdot)$ is treated as a random variable whose mean, variance and covariance with the hidden states are obtained in closed form (\S\ref{S:Phi}). The factors are then multiplied with the Gaussian multiplicative approximation (GMA), so that the class probability $Z^{\Pr}_{\mathcal C}$ is itself Gaussian with a covariance with respect to the output layer (\S\ref{S:HSM}).
\item The fixed scale $1/\alpha$ is replaced by a Gaussian \emph{gain} $G_h$ at each node, learned from the class labels through an auxiliary observation model (\S\ref{S:cal}). The network keeps learning in logit space from $\pm1$ targets exactly as in TAGI; the auxiliary channel updates the gains only.
\end{enumerate}
Section~\ref{S:ex} illustrates everything on the two-class, single-node case. Notation: $\varphi$ is the standard normal PDF, $\Phi_2(\cdot,\cdot;\rho)$ the standard bivariate normal CDF with correlation $\rho$, and
\begin{equation}
T(h,b)=\frac{1}{2\pi}\int_0^{b}\frac{e^{-h^2(1+t^2)/2}}{1+t^2}\,\mathrm dt
\end{equation}
is Owen's $T$ function (\texttt{scipy.special.owens\_t}). As everywhere in TAGI, the hidden states of a same layer are independent.

\section{Moments of $\Phi(X)$ for a Gaussian $X$}\label{S:Phi}
Let $X\sim\mathcal N(\mu,\sigma^2)$ and define
\begin{equation}
a=\frac{\mu}{\sqrt{1+\sigma^2}},\qquad \rho=\frac{\sigma^2}{1+\sigma^2}.
\end{equation}
Let $W,W_1,W_2\sim\mathcal N(0,1)$ be independent of $X$ and of each other, so that $\Phi(X)=\Pr(W\le X\mid X)$.

\paragraph{Mean.} $\mathbb E[\Phi(X)]=\Pr(W-X\le0)$ with $W-X\sim\mathcal N(-\mu,\,1+\sigma^2)$, hence
\begin{equation}
\mathbb E[\Phi(X)]=\Phi(a).
\label{E:mean}
\end{equation}

\paragraph{Variance.} Conditionally on $X$ the events $\{W_1\le X\}$ and $\{W_2\le X\}$ are independent, so $\mathbb E[\Phi(X)^2]=\Pr(W_1-X\le0,\,W_2-X\le0)=\Phi_2(a,a;\rho)$, the pair $(W_1-X,W_2-X)$ having variances $1+\sigma^2$ and covariance $\sigma^2$. With the equal-argument identity $\Phi_2(h,h;\rho)=\Phi(h)-2T\big(h,\sqrt{(1-\rho)/(1+\rho)}\big)$ and $\sqrt{(1-\rho)/(1+\rho)}=1/\sqrt{1+2\sigma^2}$,
\begin{equation}
\operatorname{Var}[\Phi(X)]=\Phi(a)\big(1-\Phi(a)\big)-2\,T\!\Big(a,\tfrac{1}{\sqrt{1+2\sigma^2}}\Big),
\qquad
\mathbb E\big[\Phi(X)\big(1-\Phi(X)\big)\big]=2\,T\!\Big(a,\tfrac{1}{\sqrt{1+2\sigma^2}}\Big).
\label{E:var}
\end{equation}
The second expression is the expected Bernoulli variance of the indicator $\mathbb 1\{W\le X\}$ given $X$; it is used in \S\ref{S:cal}. Because $\Phi(a)(1-\Phi(a))=2T(a,1)$, the variance is also
\begin{equation}
\operatorname{Var}[\Phi(X)]=\frac{1}{\pi}\int_{b}^{1}\frac{e^{-a^2(1+t^2)/2}}{1+t^2}\,\mathrm dt,\qquad b=\frac{1}{\sqrt{1+2\sigma^2}},
\label{E:varint}
\end{equation}
which is positive by construction and avoids the cancellation between the two terms of \eqref{E:var} when $\sigma$ is small; an 8-point Gauss--Legendre rule on $[b,1]$ evaluates it to $10^{-12}$ over $|a|\le8$, so Owen's $T$ never has to be coded.

\paragraph{Covariance.} Stein's identity $\mathbb E[(X-\mu)g(X)]=\sigma^2\mathbb E[g'(X)]$ with $g=\Phi$, together with $\mathbb E[\varphi(X)]=f_{X-W}(0)=\varphi(a)/\sqrt{1+\sigma^2}$, gives
\begin{equation}
\operatorname{Cov}(X,\Phi(X))=\sigma^2\,\mathbb E[\varphi(X)]=\frac{\sigma^2\,\varphi(a)}{\sqrt{1+\sigma^2}}=\sigma^2\,\frac{\partial\,\mathbb E[\Phi(X)]}{\partial\mu}.
\label{E:cov}
\end{equation}

\section{Hierarchical softmax with a Gaussian gain}\label{S:HSM}
\subsection{Node level}
At node $h$, the gain $G_h\sim\mathcal N(\mu_{G_h},\sigma^2_{G_h})$ is independent of $Z_h$ (it is a parameter, like a weight), and the branch probability is $P_h=\Phi(s_hS_h)$ with the scaled logit $S_h=G_hZ_h$. The product is treated with the GMA: $S_h$ is taken Gaussian with its exact moments
\begin{equation}
\mu_{S_h}=\mu_{G_h}\mu_h,\qquad
\sigma^2_{S_h}=\sigma^2_{G_h}\sigma_h^2+\sigma^2_{G_h}\mu_h^2+\mu^2_{G_h}\sigma_h^2 .
\label{E:S}
\end{equation}
Applying \S\ref{S:Phi} to $X=s_hS_h$ with
\begin{equation}
a_h=\frac{s_h\,\mu_{S_h}}{\sqrt{1+\sigma^2_{S_h}}},\qquad b_h=\frac{1}{\sqrt{1+2\sigma^2_{S_h}}},
\label{E:a}
\end{equation}
gives the mean and variance of the branch probability,
\begin{equation}
\mu_{P_h}=\Phi(a_h),\qquad
\sigma^2_{P_h}=\Phi(a_h)\big(1-\Phi(a_h)\big)-2T(a_h,b_h)=\frac1\pi\int_{b_h}^{1}\frac{e^{-a_h^2(1+t^2)/2}}{1+t^2}\,\mathrm dt .
\label{E:cond}
\end{equation}
For the covariances, Stein's identity applies to each Gaussian input separately. Because $P_h$ is a smooth function of $Z_h$ and of $G_h$, and the derivative of $\mathbb E[\Phi(s_hG_hZ_h)]$ with respect to the mean of one input equals the expectation of the partial derivative with respect to that input,
\begin{equation}
\operatorname{Cov}(Z_h,P_h)=\sigma_h^2\,\frac{\partial\mu_{P_h}}{\partial\mu_h},\qquad
\operatorname{Cov}(G_h,P_h)=\sigma^2_{G_h}\,\frac{\partial\mu_{P_h}}{\partial\mu_{G_h}}
\label{E:stein}
\end{equation}
hold exactly. Differentiating $\mu_{P_h}=\Phi(a_h)$, in which $\sigma^2_{S_h}$ depends on both means, gives the closed forms
\begin{equation}
\boxed{\;
\begin{aligned}
\mu_{P_h}&=\Phi(a_h),\\
\operatorname{Cov}(Z_h,P_h)&=\frac{s_h\,\sigma_h^2\,\mu_{G_h}\,\big(1+\sigma^2_{S_h}-\sigma^2_{G_h}\mu_h^2\big)}{(1+\sigma^2_{S_h})^{3/2}}\,\varphi(a_h),\\
\operatorname{Cov}(G_h,P_h)&=\frac{s_h\,\sigma^2_{G_h}\,\mu_h\,\big(1+\sigma^2_{S_h}-\mu^2_{G_h}\sigma_h^2\big)}{(1+\sigma^2_{S_h})^{3/2}}\,\varphi(a_h).
\end{aligned}\;}
\label{E:node}
\end{equation}
For a fixed gain, $\sigma_{G_h}=0$, the mean reduces to the branch probability conditional on $G_h=g$,
\begin{equation}
m_h(g)=\Phi\!\left(\frac{s_h\,g\,\mu_h}{\sqrt{1+g^2\sigma_h^2}}\right),\qquad
m_h'(g)=\frac{s_h\,\mu_h\,\varphi\big(s_hg\mu_h/\sqrt{1+g^2\sigma_h^2}\big)}{(1+g^2\sigma_h^2)^{3/2}},
\label{E:mg}
\end{equation}
and $\operatorname{Cov}(G_h,P_h)\to\sigma^2_{G_h}m_h'(\mu_{G_h})$; both are used in \S\ref{S:cal}. The only approximation in \eqref{E:node} is the Gaussian treatment of $S_h$. It is exact when either $\sigma_h=0$ or $\sigma_{G_h}=0$, since the product is then Gaussian, and the error in between is small: for the worked example of \S\ref{S:ex} ($\mu_h=0.5$, $\sigma_h=0.5$, $G_h\sim\mathcal N(1.5,0.3^2)$) the four moments $\mu_{P_h}$, $\sigma_{P_h}$, $\operatorname{Cov}(Z_h,P_h)$, $\operatorname{Cov}(G_h,P_h)$ are $0.7229$, $0.2142$, $0.0977$, $0.00773$ against Monte Carlo $0.7226$, $0.2076$, $0.0982$, $0.00786$, and with a wide prior on the gain ($\sigma_{G_h}=1$) $\operatorname{Cov}(G_h,P_h)$ stays within $10\%$ of the exact value.

Three remarks. First, projecting through $S_h$ in the usual TAGI manner,
\begin{equation}
\operatorname{Cov}(G_h,P_h)\approx\frac{\operatorname{Cov}(G_h,S_h)\operatorname{Cov}(S_h,P_h)}{\sigma^2_{S_h}}=\frac{s_h\,\sigma^2_{G_h}\,\mu_h\,\varphi(a_h)}{\sqrt{1+\sigma^2_{S_h}}},
\end{equation}
is not equivalent to \eqref{E:node}: it lacks the factor $(1+\sigma^2_{S_h}-\mu^2_{G_h}\sigma_h^2)/(1+\sigma^2_{S_h})$, which tends to $1/(1+\mu^2_{G_h}\sigma_h^2)$ as $\sigma_{G_h}\to0$, because it ignores that a larger gain also widens $S_h$. The projection therefore exceeds \eqref{E:node} by $(1+\sigma^2_{S_h})/(1+\sigma^2_{S_h}-\mu^2_{G_h}\sigma_h^2)\approx1+\mu^2_{G_h}\sigma_h^2$: $1.54$ for the worked example of \S\ref{S:ex}, $3.0$ at $\sigma_h=1$ with $\mu_{G_h}=1.5$, and between $1.4$ and $6.7$ over the stream of \S\ref{S:ex} once the gain is near $2$. Used in \S\ref{S:cal} it makes the gain posterior over-confident; the Stein form is the one to use. Second, the gain update of \S\ref{S:cal} requires $\operatorname{Cov}(G_h,P_h)^2\le\sigma^2_{G_h}\mu_{P_h}(1-\mu_{P_h})$. With exact moments this is Cauchy--Schwarz; with \eqref{E:node} it is not guaranteed, but over $2\times10^6$ random combinations covering $|\mu_h|\le8$, $\sigma_h\le4$, $-1\le\mu_{G_h}\le6$ and $\sigma_{G_h}\le5$ the ratio never exceeded $0.64$ (a floor at zero is kept anyway, \S\ref{S:cal}). Third, $\operatorname{sign}\operatorname{Cov}(G_h,P_h)=\operatorname{sign}(s_h\mu_h)$, since $1+\sigma^2_{S_h}-\mu^2_{G_h}\sigma_h^2=1+\sigma^2_{G_h}(\sigma_h^2+\mu_h^2)>0$. The sibling branch at the same node is $1-P_h$: same variance, mean $1-\mu_{P_h}$, covariances of opposite sign. With the gain fixed at $1/\alpha$, $\mu_{P_h}$ is exactly one factor of \eqref{E:tagi}.

\subsection{Class level}
The class probability $Z^{\Pr}_{\mathcal C}=\prod_{h=1}^{\mathtt H}P_h$ is built recursively as $Q_1=P_1$ and $Q_h=Q_{h-1}P_h$, using the GMA with $\operatorname{Cov}(Q_{h-1},P_h)=0$ (independent nodes). Writing $X_l$ for either $Z_l$ or $G_l$,
\begin{align}
\mu_{Q_h}&=\mu_{Q_{h-1}}\mu_{P_h},\\
\sigma^2_{Q_h}&=\sigma^2_{Q_{h-1}}\sigma^2_{P_h}+\mu^2_{Q_{h-1}}\sigma^2_{P_h}+\mu^2_{P_h}\sigma^2_{Q_{h-1}},\\
\operatorname{Cov}(Q_h,X_l)&=\begin{cases}\mu_{P_h}\operatorname{Cov}(Q_{h-1},X_l)& l<h,\\ \mu_{Q_{h-1}}\operatorname{Cov}(P_h,X_h)& l=h.\end{cases}
\end{align}
Unrolled, for $l$ on the path,
\begin{equation}
\boxed{\;
\begin{aligned}
\mu_{\Pr}&=\prod_{h=1}^{\mathtt H}\mu_{P_h},\\
\sigma^2_{\Pr}&=\prod_{h=1}^{\mathtt H}\big(\sigma^2_{P_h}+\mu_{P_h}^2\big)-\prod_{h=1}^{\mathtt H}\mu_{P_h}^2,\\
\operatorname{Cov}(Z^{\Pr}_{\mathcal C},X_l)&=\operatorname{Cov}(P_l,X_l)\prod_{h\neq l}\mu_{P_h},
\end{aligned}\;}
\label{E:class}
\end{equation}
and $\operatorname{Cov}(Z^{\Pr}_{\mathcal C},X_l)=0$ for output units off the path. With independent nodes these are the exact moments of the product given the node moments; the Gaussian assumption only enters when $Z^{\Pr}_{\mathcal C}$ is subsequently used as a Gaussian. Since $\mu_{P_h}+(1-\mu_{P_h})=1$ at every node, $\sum_{\mathcal C}\mu_{\Pr}=1$ exactly when $\mathtt K$ is a power of two (otherwise normalize as in TAGI).

\section{Learning the gains from the labels}\label{S:cal}
\subsection{Which channel learns what}
The network is trained as in TAGI: for a training sample of class $\mathcal C^\ast$, the $\mathtt H$ path units receive the logit observations $y_h=z_h+v$ with targets $y_h=s_h\in\{-1,+1\}$ and $V\sim\mathcal N(0,\sigma_V^2)$, and the backward pass updates hidden states and parameters. These $\pm1$ targets pin the scale of $Z_h$; the gain $G_h$ therefore only decides how much probability mass one unit of logit is worth, i.e.\ it is a learned per-node calibration (a Platt/temperature scaling). Without the logit channel, $G_h$ would be confounded with the scale of the output weights.

The gains are learned from labels through a second, \emph{auxiliary} observation model in probability space, and the natural place to run it is the validation split, once per epoch, sample by sample: the network is not updated there, so the forward-pass moments $(\mu_h,\sigma_h^2)$ are genuine out-of-sample predictions, and the recursion is a scalar per node, so the pass costs $N_{\text{val}}\,\mathtt H$ evaluations of $\Phi$ and $\varphi$ and needs no batching. This is temperature scaling fitted on the validation set, inside TAGI, after every epoch. The training loop is left exactly as it is.

If the channel is instead run on the training stream (e.g.\ online learning without epochs), two rules apply. It must not update $\bm Z$ (the label is already used by the logit channel), and it must use the forward-pass moments rather than the moments after the logit update: calibration relates a \emph{prediction} to an outcome, the post-update moments are a blend of the two, and a channel fed with them sees predictions that already agree with the label and sharpens the gain without bound (\S\ref{S:ex}). Under these two rules the two channels commute. Training-set predictions are out-of-sample only during the first pass over the data, so the first epoch is the only one that yields an honest calibration on the training stream.

\subsection{Auxiliary observation model}
The categorical likelihood of the label factorizes over the path, $Z^{\Pr}_{\mathcal C^\ast}=\prod_hP_h$, so observing the class is equivalent to observing ``branch taken'' at each node: $y_h=1$ with $y_h\mid P_h\sim\mathrm{Bernoulli}(P_h)$. The observation is always $1$ because the direction is carried by the sign: $P_h=\Phi(s_hG_hZ_h)$ with $s_h=(-1)^{\mathcal C^\ast_h}$ is, by construction, the probability of the branch the class took. Equivalently one may fix the orientation, $P_h^{+}=\Phi(G_hZ_h)$, and observe $y_h^{+}=1-\mathcal C^\ast_h\in\{0,1\}$; since $\mu_{P_h^{-}}=1-\mu_{P_h^{+}}$ and $\operatorname{Cov}(G_h,P_h^{-})=-\operatorname{Cov}(G_h,P_h^{+})$, the update for $y_h^{+}=0$ coincides term by term with the $s_h$-oriented update for $y_h=1$, so the two conventions are the same computation. In Gaussian form,
\begin{equation}
Y_h=P_h+V_h,\qquad V_h\sim\mathcal N(0,\sigma^2_{V_h}),\qquad
\sigma^2_{V_h}=\mathbb E[P_h(1-P_h)]=\mu_{P_h}(1-\mu_{P_h})-\sigma^2_{P_h}.
\label{E:obs}
\end{equation}
The choice of $\sigma^2_{V_h}$ follows from the law of total variance,
\begin{equation}
\operatorname{Var}[Y_h]=\mathbb E\big[\operatorname{Var}[Y_h\mid P_h]\big]+\operatorname{Var}\big[\mathbb E[Y_h\mid P_h]\big]=\mathbb E[P_h(1-P_h)]+\sigma^2_{P_h}=\mu_{P_h}(1-\mu_{P_h}),
\end{equation}
the Bernoulli variance. Taking $\sigma^2_{V_h}=\mu_{P_h}(1-\mu_{P_h})$ instead would count $\sigma^2_{P_h}$ twice and damp every update by the factor $\mu_{P_h}(1-\mu_{P_h})/\big(\mu_{P_h}(1-\mu_{P_h})+\sigma^2_{P_h}\big)$. In practice only $\operatorname{Var}[Y_h]=\mu_{P_h}(1-\mu_{P_h})$ is needed, so $\sigma^2_{P_h}$ is never evaluated for the update.

The auxiliary channel is the two-variable joint Gaussian
\begin{equation}
\begin{bmatrix}G_h\\ Y_h\end{bmatrix}\sim\mathcal N\!\left(\begin{bmatrix}\mu_{G_h}\\ \mu_{P_h}\end{bmatrix},
\begin{bmatrix}\sigma^2_{G_h}&\operatorname{Cov}(G_h,P_h)\\ \operatorname{Cov}(G_h,P_h)&\mu_{P_h}(1-\mu_{P_h})\end{bmatrix}\right),
\end{equation}
and conditioning on $y_h=1$ gives the gain update
\begin{equation}
\boxed{\;
\mu_{G_h\mid y}=\mu_{G_h}+\frac{\operatorname{Cov}(G_h,P_h)}{\mu_{P_h}},\qquad
\sigma^2_{G_h\mid y}=\sigma^2_{G_h}-\frac{\operatorname{Cov}(G_h,P_h)^2}{\mu_{P_h}(1-\mu_{P_h})}\;}
\label{E:update}
\end{equation}
with $\mu_{P_h}$ and $\operatorname{Cov}(G_h,P_h)$ from \eqref{E:node}. $Z_h$ enters only through $(\mu_h,\sigma_h^2)$ inside these two quantities; it is not a variable of the joint and receives no update. The bound discussed in \S\ref{S:HSM} is what keeps $\sigma^2_{G_h\mid y}\ge0$.

Three properties. First, the mean in \eqref{E:update} is exact under the Bernoulli likelihood, without any Gaussian assumption on $P_h$: $\mathbb E[G_h\mid y_h=1]=\mathbb E[G_hP_h]/\mathbb E[P_h]=\mu_{G_h}+\operatorname{Cov}(G_h,P_h)/\mu_{P_h}$, and likewise $\mathbb E[G_h\mid y_h=0]=\mu_{G_h}-\operatorname{Cov}(G_h,P_h)/(1-\mu_{P_h})$, which is the same linear form evaluated at $y_h=0$. Second, the variance step is exact \emph{on average over the label}: for a binary $Y_h$, $\mathbb E[G_h\mid Y_h]$ is affine in $Y_h$, so the law of total variance gives $\sigma^2_{G_h}-\operatorname{Cov}(G_h,P_h)^2/\operatorname{Var}[Y_h]=\mathbb E_{y_h}\big[\operatorname{Var}[G_h\mid y_h]\big]$ exactly. What it discards is the label dependence. For $\sigma_h$ of practical size, where $m_h(g)$ saturates in $g$, the exact posterior is narrower after a confirming label and wider after a surprising one ($0.2984$ and $0.3022$ for the worked example of \S\ref{S:ex}, around the recursion's $0.2995$); in the limit $\sigma_h\to0$, where $m_h$ is a plain probit in $g$, the order reverses ($0.2985$ and $0.2975$), the crossover being near $\sigma_h\approx0.15$ for that example. In either case the recursion only ever shrinks, and over a sequence it ends somewhat over-confident (quantified in \S\ref{S:ex}); this is what the variance re-initialization and, for streams, the process noise of \S\ref{S:remarks} compensate. Third, $\operatorname{sign}\operatorname{Cov}(G_h,P_h)=\operatorname{sign}(s_h\mu_h)$: the gain increases when the network already agrees with the branch taken (sharpen) and decreases when it disagrees (soften), which is the behaviour expected from a calibration parameter. With the small-$\sigma_{G_h}$ form $\operatorname{Cov}(G_h,P_h)\approx\sigma^2_{G_h}m_h'(\mu_{G_h})$ of \eqref{E:mg}, the mean step reads $\sigma^2_{G_h}\,\partial_g\log\mathrm{Bern}\big(y_h;m_h(g)\big)$ and the variance step adds the Fisher information: the recursion is online Fisher scoring for the maximum-likelihood gain, i.e.\ temperature scaling fitted online, per node.

\paragraph{Why not observe the product.} One could instead observe $y=1$ on $Z^{\Pr}_{\mathcal C^\ast}$ with $\operatorname{Var}[Y]=\mu_{\Pr}(1-\mu_{\Pr})$. By \eqref{E:class} the resulting mean update of $G_l$ is $\operatorname{Cov}(Z^{\Pr}_{\mathcal C^\ast},G_l)/\mu_{\Pr}=\operatorname{Cov}(P_l,G_l)/\mu_{P_l}$, identical to \eqref{E:update}. The variance updates differ: the product version removes $\operatorname{Cov}(P_l,G_l)^2\,\mu_{\Pr}/\big(\mu_{P_l}^2(1-\mu_{\Pr})\big)$, i.e.\ a fraction $\mu_{\Pr}(1-\mu_{P_l})/\big(\mu_{P_l}(1-\mu_{\Pr})\big)\le1$ of the node version, because the exact posterior factorizes over nodes and the product formulation loses that structure. The node-wise form \eqref{E:update} is thus both simpler and more accurate; the product \eqref{E:class} is only needed to report class probabilities. Likewise, pushing $y=0$ on the other classes adds nothing: at a path node the sibling's ``0'' is the same event as the taken branch's ``1'', and the likelihood of the label does not involve the gains of nodes off the path.

\paragraph{Sharing gains.} Nothing requires one gain per node. A training sample visits exactly one node per level of the tree, so gains shared across all nodes of a level ($\mathtt H$ gains instead of $\mathtt K-1$) receive exactly one update per sample with the algorithm unchanged; this pools the data of deep nodes, which individually see only the samples of their subtree. A single global gain (one parameter, the closest analogue of temperature scaling) is also possible: the $\mathtt H$ updates of a sample are then applied sequentially, each with $\mu_{P_h}$ and $\operatorname{Cov}(G,P_h)$ recomputed from the gain moments left by the previous node. With a validation split of $N_{\text{val}}$ samples a node at the deepest level sees only about $2N_{\text{val}}/\mathtt K$ of them, so per-level sharing is the sensible default; per-node gains are warranted only when the tree is shallow or the validation set large.

\paragraph{Epochs.} Unlike the weights, the gains are not learned across epochs. A gain calibrates the forward-pass predictions of the network \emph{as it currently is}; its posterior is conditional on $\bm\theta$, and the precision accumulated during epoch $e$ describes the calibration of a network that no longer exists once epoch $e+1$ has started. Formally the gain of successive epochs follows a random walk, $G_h^{(e+1)}=G_h^{(e)}+W$, $W\sim\mathcal N(0,\sigma^2_{G,0})$, so the prior at the start of an epoch keeps the previous posterior mean and has variance $\sigma^2_{G_h}+\sigma^2_{G,0}\approx\sigma^2_{G,0}$: keep $\mu_{G_h}$, re-initialize $\sigma^2_{G_h}\leftarrow\sigma^2_{G,0}$ before each validation pass. Since a gain converges within a few hundred samples (Fig.~\ref{F:cal}a) a validation split of a few thousand samples is ample time to re-converge, and because the channel is decoupled from the network the reset costs nothing in training. On a stream whose calibrated gain drifts from $2.0$ to $1.0$ over five epochs (as logits sharpen during training), the recursion without reset ends at $1.19\pm0.04$, confidently wrong; with the reset it ends at $1.00\pm0.05$; a per-sample process noise $q=10^{-4}$ also tracks but with $\sigma_G\approx0.2$, and remains the tool for streams without epochs.

\paragraph{Numerical guards.} When $\Phi$ saturates, $\mu_{P_h}(1-\mu_{P_h})$ underflows and \eqref{E:update} divides by zero; in float32 this happens for a fraction of a percent of the samples once the gain is around $2$. Clip $\mu_{P_h}$ to $[\epsilon,1-\epsilon]$ with $\epsilon\approx10^{-6}$ before \eqref{E:update}, and floor $\sigma^2_{G_h\mid y}$ at zero, since the bound of \S\ref{S:HSM} is not guaranteed for the GMA moments and roundoff can violate it in any case.

\subsection{Algorithm}
\begin{algorithm}[H]
\caption{Gain calibration pass, run on the validation split after each epoch (gains initialized once with a weakly informative prior, e.g.\ $\mu_{G_h}=0.3$, $\sigma_{G,0}=1$). The training loop is standard TAGI and is not modified.}
Keep $\mu_{G_h}$, re-initialize $\sigma^2_{G_h}\leftarrow\sigma^2_{G,0}$ for all gains\;
Forward pass on the validation split (no network update)\;
\ForEach{validation sample of class $\mathcal C^\ast$, in sequence}{
\For{$h=1$ \KwTo $\mathtt H$}{
Take the forward-pass $(\mu_h,\sigma_h^2)$ of the path unit $h$ and $s_h=(-1)^{\mathcal C^\ast_h}$\;
Compute $\sigma^2_{S_h}$, $a_h$ by \eqref{E:S}--\eqref{E:a}, then $\mu_{P_h}=\Phi(a_h)$ and $\operatorname{Cov}(G_h,P_h)$ by \eqref{E:node}\;
Clip $\mu_{P_h}$ to $[\epsilon,1-\epsilon]$; update $(\mu_{G_h},\sigma^2_{G_h})$ with \eqref{E:update}, using $y_h=1$, and floor $\sigma^2_{G_h}$ at zero\;
}
}
Report class probabilities and their variances with \eqref{E:class} using the updated gains\;
\end{algorithm}

\subsection{Prediction}
At test time the class probabilities are $\mu_{\Pr}$ from \eqref{E:class} with the current $(\mu_{G_h},\sigma^2_{G_h})$; $\sigma^2_{\Pr}$ quantifies how uncertain the probability itself is, and $\operatorname{Cov}(Z^{\Pr}_{\mathcal C},Z_l)$ is available for any downstream inference on the output layer.

\section{Illustration: two classes, one node}\label{S:ex}
With $\mathtt K=2$ the tree has a single node, $\mathtt H=1$: one output unit $Z\sim\mathcal N(\mu_Z,\sigma_Z^2)$, one gain $G\sim\mathcal N(\mu_G,\sigma_G^2)$, $s=+1$ for class~1 and $s=-1$ for class~0. The probability of class~1 is $P=\Phi(GZ)$ with moments \eqref{E:node} at $s=1$; class~0 has probability $1-P$.

\paragraph{Node moments.} Take $\mu_Z=0.5$, $\sigma_Z=0.5$ and $G\sim\mathcal N(1.5,0.3^2)$. Equations \eqref{E:cond}--\eqref{E:node} give $\mu_P=0.723$, $\sigma_P=0.214$, $\operatorname{Cov}(Z,P)=0.098$ and $\operatorname{Cov}(G,P)=7.7\times10^{-3}$; a Monte Carlo estimate with $2\times10^6$ samples returns $0.723$, $0.208$, $0.098$ and $7.9\times10^{-3}$. Figure~\ref{F:node}a compares the actual distribution of $P$ with the moment-matched Gaussian: the shape is skewed, but only the mean, variance and covariances enter the linear update. Figure~\ref{F:node}b shows how $\sigma_Z$ flattens $\mu_P$ and creates the spread $\sigma_P$ that the expected-value formulation \eqref{E:tagi} does not report. Figure~\ref{F:node}c shows the two covariances as functions of $\mu_Z$: $\operatorname{Cov}(Z,P)$ is a positive bump centred at the decision boundary, whereas $\operatorname{Cov}(G,P)$ is odd in $\mu_Z$ and changes sign with the side of the boundary on which the prediction lies.

\begin{figure}[h]
\centering\includegraphics[width=\textwidth]{fig_node.pdf}
\caption{Node level for $\mathtt K=2$. (a) Distribution of $P=\Phi(GZ)$ and its moment-matched Gaussian for the worked example; the dotted line is $\mu_P$. (b) $\mu_P\pm\sigma_P$ as a function of the forward-pass mean $\mu_Z$ for three values of $\sigma_Z$, at fixed $G=1.5$. (c) $\operatorname{Cov}(Z,P)$ and $\operatorname{Cov}(G,P)/\sigma_G^2$ as functions of $\mu_Z$ for $\sigma_Z=1$, $G\sim\mathcal N(1.5,0.3^2)$.}
\label{F:node}
\end{figure}

\paragraph{One observation.} Suppose the label is class~1. The gain update \eqref{E:update} gives $\Delta\mu_G=\operatorname{Cov}(G,P)/\mu_P=+0.0107$ and $\sigma_G\to0.2995$; had the label been class~0, $\Delta\mu_G=-\operatorname{Cov}(G,P)/(1-\mu_P)=-0.0279$. The exact posterior means, computed on a grid over $g$, are $+0.0109$ and $-0.0284$. The asymmetry is the Bernoulli likelihood at work: the prediction favoured class~1 with probability $0.72$, so confirming it is weak evidence and contradicting it is strong evidence. Figure~\ref{F:update}a plots $\Delta\mu_G$ against the forward-pass mean for both labels: a prediction on the correct side of the boundary produces a small step up (sharpen), a prediction on the wrong side a large step down (soften); the two curves are mirror images, as they must be since the two labels are complementary events. Figures~\ref{F:update}b--c show one sample ($\mu_Z=0.8$, $\sigma_Z=0.8$, $G\sim\mathcal N(1,0.3^2)$, $\sigma_V=0.3$) as it is seen by the two channels: in logit space the target $s=+1$ moves and narrows $Z$ (and, through the backward pass, $\bm\theta$); in probability space the label $y_h=1$ is compared with the forward-pass distribution of $P$ and moves $G$ only.

\begin{figure}[h]
\centering\includegraphics[width=\textwidth]{fig_update.pdf}
\caption{One training sample. (a) Mean update of the gain as a function of the forward-pass mean $\mu_Z$, for a label equal to class~1 (blue) or class~0 (red); $\sigma_Z=1$, $G\sim\mathcal N(1,0.3^2)$. (b) Logit channel: forward-pass density of $Z$, target $s=+1$, and the density after the logit update. (c) Probability channel for the same sample: distribution of $P=\Phi(GZ)$ obtained from the forward-pass moments, its moment-matched Gaussian, and the label $y_h=1$.}
\label{F:update}
\end{figure}

\paragraph{Learning the gain from a validation pass.} A set of $N=3000$ samples, standing for a validation split, is generated as follows. The network's forward-pass moments are $\mu_i\sim\mathcal N(0,1.5^2)$ and $\sigma_i\sim\mathcal U(0.3,1.2)$; the true logit is $z_i\sim\mathcal N(\mu_i,\sigma_i^2)$, i.e.\ the network is calibrated in logit space; the label is $y_i\sim\mathrm{Bernoulli}\big(\Phi(g^\ast z_i)\big)$ with $g^\ast=2$. The calibrated gain is therefore $g^\ast$, and the initial value $\mu_G=0.3$ is badly under-confident. Because $G$ is a scalar, the exact posterior $p(g\mid\text{data})\propto p(g)\prod_i m_i(g)^{y_i}\big(1-m_i(g)\big)^{1-y_i}$ can be tabulated on a grid and used as a reference for the recursion. Figure~\ref{F:cal}a shows the weakly informative prior $\mathcal N(0.3,1^2)$: processed in sequence with Algorithm~1, the recursion reaches $2.02\pm0.14$ after 3000 samples, the exact posterior $2.14\pm0.20$. Over four independent streams the recursion mean stays within $0.15$ of the exact mean while its $\sigma_G$ is $70$--$95\%$ of the exact width, the over-confidence anticipated in \S\ref{S:cal}. Figure~\ref{F:cal}b shows two things to avoid. With a tight prior $\mathcal N(0.3,0.3^2)$ the recursion ends at $1.49\pm0.07$ against an exact posterior of $1.77\pm0.12$: the prior itself is responsible for part of the lag (it places the truth $5.7\sigma$ away, and exact Bayes has not recovered after 3000 samples either), and the Gaussian recursion adds a further $0.3$ with a $\sigma_G$ at $60\%$ of the exact width; hence the recommendation of a wide prior and of a process noise $q$ whose natural target is this gap. Feeding the channel the moments \emph{after} the logit update (\S\ref{S:cal}) makes $G$ grow without bound ($3.7$ and rising), because the channel then sees predictions that already agree with the label. Figure~\ref{F:cal}c is the reliability diagram on $40\,000$ fresh test samples: with $G=0.3$ the probabilities are under-confident (ECE $0.197$, NLL $0.511$); with the learned gain they lie on the diagonal (ECE $0.005$, NLL $0.363$), identical to what the true gain achieves.

\begin{figure}[h]
\centering\includegraphics[width=\textwidth]{fig_calibration.pdf}
\caption{Online calibration for $\mathtt K=2$. (a) Prior $\mathcal N(0.3,1^2)$: recursion \eqref{E:update} (blue, with the $\pm2\sigma_G$ band) against the exact posterior of $G$ on the same data (red, mean $\pm2\sigma$); dotted: true gain. (b) Same with the tight prior $\mathcal N(0.3,0.3^2)$; orange dashed: channel fed with post-update moments (wide prior). (c) Reliability diagram on independent test data with the initial and the learned gain.}
\label{F:cal}
\end{figure}

\section{Remarks}\label{S:remarks}
\begin{itemize}
\item \textbf{Limits.} For $\sigma_h\to0$, $\operatorname{Cov}(Z_h,P_h)\to0$ but $\sigma^2_{P_h}$ and $\operatorname{Cov}(G_h,P_h)$ tend to the finite (and, since $S_h=\mu_hG_h$ is then Gaussian, exact) limits $\operatorname{Var}[\Phi(s_hG_h\mu_h)]$ and $\operatorname{Cov}(G_h,\Phi(s_hG_h\mu_h))$, equal to $1.8\times10^{-3}$ and $0.012$ for $\mu_h=1$ and $G_h\sim\mathcal N(1.5,0.3^2)$: a confident logit still calibrates the gain. For $\mu_h=0$, $\sigma^2_{P_h}=\tfrac14-\tfrac1\pi\arctan b_h$, which tends to $1/4$ as $\sigma_{S_h}\to\infty$ (a fair coin). With a fixed gain $1/\alpha$, $\mu_{\Pr}$ is exactly TAGI's Eq.~\eqref{E:tagi}.
\item \textbf{Prior on the gain.} Use a weakly informative prior, $\sigma_{G,0}\gtrsim\mu_{G_h}$, both at initialization and at each epoch reset. A tight prior slows down exact Bayes and, on top of that, makes the Gaussian recursion over-confident before the mean has converged (Fig.~\ref{F:cal}b).
\item \textbf{Adaptivity.} Within a pass $\sigma^2_{G_h}$ only decreases, and the recursion is in any case narrower than the exact posterior (Fig.~\ref{F:cal}a). Across epochs the variance is re-initialized before each validation pass (\S\ref{S:cal}), which is what keeps the gains following the network; a per-sample process noise $q$ does the same for streams without epochs, and the gap to the exact width in \S\ref{S:ex} gives $q$ a target.
\item \textbf{Positivity.} Nothing enforces $G_h>0$; a sign flip would silently swap the two branches of a node, so $\mu_{G_h}$ should be clipped at a small positive value.
\item \textbf{Cost.} The calibration pass needs one $\Phi$ and one $\varphi$ per path node per validation sample, as a scalar recursion per node; the recursion is sequential by nature (each step recomputes the moments from the updated gain), which is why it is run sample by sample rather than in batches --- a batch form exists (linear regression of $P_h$ on $G_h$, additive in information form) but adds an approximation that degrades with the batch size, and there is nothing to gain from it here. The prediction path additionally needs $\sigma^2_{P_h}$, computed with the integral form in \eqref{E:cond} and an 8-point Gauss--Legendre rule on $[b_h,1]$: branch-free, positive by construction, accurate to float32 roundoff, and free of Owen's $T$.
\item \textbf{Verification.} All closed forms were checked against Monte Carlo and against the exact node moments obtained by conditioning on $G_h$ and integrating with a 20-point Gauss--Hermite rule: the node moments \eqref{E:node}, the class moments \eqref{E:class}, and the exact posterior means $\mathbb E[G_hP_h]/\mathbb E[P_h]$ and $\mathbb E[G_h(1-P_h)]/\mathbb E[1-P_h]$ against \eqref{E:update} (\S\ref{S:ex}).
\end{itemize}

\end{document}
